% © 2026 Ing. David Jaroš — CC BY-NC-ND 4.0
%
% This work is licensed under a Creative Commons Attribution-NonCommercial-NoDerivatives
% 4.0 International License (CC BY-NC-ND 4.0).
%
% License History: Earlier drafts (up to v0.3) were released under CC BY 4.0.
% From v0.4 onward, all material is released under CC BY-NC-ND 4.0 to protect
% the integrity of the theoretical work during ongoing academic development.
%
% See LICENSE.md for full license text.

\ifdefined\INCLUDEMODE
\else
  \documentclass[12pt,a4paper]{article}
  \usepackage[a4paper, margin=2.5cm]{geometry}
  \usepackage{amsmath, amssymb, amsthm}
  \usepackage{mathtools}
  \usepackage{hyperref}
  \usepackage{xcolor}
  \usepackage{mdframed}
  \usepackage{booktabs}

  \newtheorem{theorem}{Theorem}[section]
  \newtheorem{lemma}[theorem]{Lemma}
  \newtheorem{proposition}[theorem]{Proposition}
  \newtheorem{conjecture}[theorem]{Conjecture}
  \theoremstyle{definition}
  \newtheorem{definition}[theorem]{Definition}
  \theoremstyle{remark}
  \newtheorem{remark}[theorem]{Remark}

  \newmdenv[backgroundcolor=yellow!20, linecolor=orange, linewidth=1.5pt]{gapbox}
  \newmdenv[backgroundcolor=green!15, linecolor=green!60!black, linewidth=1.5pt]{resultbox}
  \newmdenv[backgroundcolor=red!10, linecolor=red!60, linewidth=1.5pt]{deadendbox}
  \newmdenv[backgroundcolor=blue!8, linecolor=blue!50, linewidth=1.5pt]{insightbox}

  \title{\textbf{Gap A --- Approach A3: One-Loop $B_{\mathrm{eff}}$ on $T^2(R_t, R_\psi)$}\\
         \large Two-Circle Kaluza-Klein Computation}
  \author{David Jaro\v{s}}
  \date{2026-03-06}

  \begin{document}
  \maketitle

  \begin{abstract}
  We extend the one-loop Kaluza-Klein (KK) computation of the effective
  coefficient $B_{\mathrm{eff}}$ from the single $\psi$-circle $S^1(R_\psi)$
  to the full complex-time torus $T^2(R_t, R_\psi)$.  On the full torus the
  KK sum depends on both winding numbers $(m, n)$ and on the modular parameter
  $\tau = i R_\psi / R_t$.  We compute $B_{\mathrm{eff}}$ as a function of the
  radius ratio $r = R_t/R_\psi$ and test whether $B_{\mathrm{eff}} = 46.3$ is
  achieved for $r = \sqrt{137}$ or $r = \sqrt{139}$.  The analysis shows
  that for $r \gg 1$ the correction is exponentially small
  ($\sim e^{-2\pi r}$), giving $B_{\mathrm{eff}} \approx B_0 = 8\pi$.
  The only regime that produces $B_{\mathrm{eff}} = 46.3$ requires $r \sim 1$,
  which implies both circles are compact at the same scale --- a scenario that
  conflicts with Lorentz invariance at accessible energies.  This approach is
  a dead end.
  \end{abstract}
\fi

%──────────────────────────────────────────────────────────────────────────────
\section{Setup: Complex Time on $T^2(R_t, R_\psi)$}
\label{sec:setup}
%──────────────────────────────────────────────────────────────────────────────

In UBT complex time is $\tau = t + i\psi$ with $t$ real and $\psi$
imaginary.  In the minimal setting $\psi$ is compactified on $S^1(R_\psi)$
with circumference $2\pi R_\psi$.  In the extended two-circle setting,
$t$ is also compactified on $S^1(R_t)$, so that
\begin{equation}
  \tau \;\in\; T^2(R_t, R_\psi)
  \;:=\; \mathbb{R}/(2\pi R_t\mathbb{Z}) \times \mathbb{R}/(2\pi R_\psi\mathbb{Z}).
  \label{eq:torus}
\end{equation}
This torus has the modular parameter
\begin{equation}
  \hat\tau \;=\; i\,\frac{R_\psi}{R_t}.
  \label{eq:modparam}
\end{equation}
The case $R_t \to \infty$ (decompactified real time) recovers the single
$\psi$-circle $S^1(R_\psi)$.

\subsection{KK spectrum on $T^2$}

A field $\Phi(\tau)$ with the double-periodicity of $T^2$ has the Fourier
expansion
\begin{equation}
  \Phi(\tau) \;=\; \sum_{m,n \in \mathbb{Z}} \Phi_{m,n}\,
  \exp\!\left(\frac{imt}{R_t} + \frac{in\psi}{R_\psi}\right),
  \label{eq:KK_modes}
\end{equation}
with KK masses
\begin{equation}
  M_{m,n}^2 \;=\; \frac{m^2}{R_t^2} + \frac{n^2}{R_\psi^2}.
  \label{eq:KK_mass}
\end{equation}

%──────────────────────────────────────────────────────────────────────────────
\section{One-Loop $B_{\mathrm{eff}}$ on $T^2$}
\label{sec:one_loop}
%──────────────────────────────────────────────────────────────────────────────

The one-loop effective coefficient $B_{\mathrm{eff}}$ is obtained by summing
the KK contribution over all winding modes.  Schematically,
\begin{equation}
  B_{\mathrm{eff}}(R_t, R_\psi)
  \;=\; B_0 \;\times\;
  \sum_{m,n \in \mathbb{Z}} \mathcal{F}(m, n;\, R_t, R_\psi),
  \label{eq:Beff_sum}
\end{equation}
where $\mathcal{F}$ is a dimensionless function encoding the contribution of
the $(m,n)$ mode.  In the limit $R_t \to \infty$, only $m = 0$ modes survive
and $\mathcal{F}(0,n;\infty, R_\psi) = \delta_{n,\pm 1}$, giving
\begin{equation}
  B_{\mathrm{eff}} \;\xrightarrow{R_t \to \infty}\;
  B_0 \;=\; 8\pi,
  \label{eq:Beff_limit}
\end{equation}
which is the proved one-circle result.

\subsection{Theta-function representation}

Using the Poisson resummation formula, the sum \eqref{eq:Beff_sum} can be
written in terms of the Jacobi theta function $\vartheta_3$:
\begin{equation}
  \sum_{m \in \mathbb{Z}} e^{-\pi m^2 R_t^2 / R_\psi^2}
  \;=\; \vartheta_3\!\left(0,\, e^{-\pi R_t^2/R_\psi^2}\right)
  \;=\; \sqrt{\frac{R_\psi}{R_t}}\;\vartheta_3\!\left(0,\, e^{-\pi R_\psi^2/R_t^2}\right).
  \label{eq:theta_resum}
\end{equation}
For the ratio $r := R_t / R_\psi \gg 1$ this gives
\begin{equation}
  B_{\mathrm{eff}}(r)
  \;=\; B_0 \left(1 + 2\,e^{-2\pi r} + O(e^{-4\pi r})\right),
  \label{eq:Beff_expansion}
\end{equation}
showing that the correction to $B_0$ is \emph{exponentially small} for large $r$.

%──────────────────────────────────────────────────────────────────────────────
\section{Test at $r = \sqrt{137}$ and $r = \sqrt{139}$}
\label{sec:test_r}
%──────────────────────────────────────────────────────────────────────────────

The UBT prime attractors suggest testing the special radii
$R_t / R_\psi = \sqrt{137}$ and $R_t / R_\psi = \sqrt{139}$.

\paragraph{$r = \sqrt{137} \approx 11.70$.}
From \eqref{eq:Beff_expansion}:
\begin{equation}
  B_{\mathrm{eff}}(\sqrt{137})
  \;=\; 8\pi\left(1 + 2\,e^{-2\pi\sqrt{137}} + \ldots\right).
\end{equation}
The correction is $2\,e^{-2\pi\times 11.70} \approx 2\,e^{-73.5}$, which is
$\approx 10^{-32}$ --- completely negligible.

\paragraph{$r = \sqrt{139} \approx 11.79$.}
Similarly, the correction is $\approx 2\,e^{-74.1} \approx 10^{-32}$.

\begin{deadendbox}
\textbf{Dead end (large-ratio limit):}
At $r = \sqrt{137}$ or $r = \sqrt{139}$, the two-circle correction to
$B_0 = 8\pi$ is $< 10^{-32}$.  The torus geometry with large $r$ reproduces
only $B_0 = 8\pi \approx 25.13$, not $B_{\mathrm{phenom}} \approx 46.3$.
\end{deadendbox}

%──────────────────────────────────────────────────────────────────────────────
\section{What Radius Ratio Would Give $B_{\mathrm{eff}} = 46.3$?}
\label{sec:required_r}
%──────────────────────────────────────────────────────────────────────────────

To achieve $B_{\mathrm{eff}} = 46.3 = 1.842 \times B_0$, one needs the full
KK sum to contribute an extra factor of $0.842$.  From the theta-function
representation \eqref{eq:theta_resum}, the ratio $B_{\mathrm{eff}}/B_0 = 1.842$
requires
\begin{equation}
  \vartheta_3\!\left(0, e^{-\pi/r^2}\right) \;=\; 1.842 \times \vartheta_3(0, 0) = 1.842,
  \label{eq:theta_condition}
\end{equation}
which, by the asymptotic expansion of $\vartheta_3$, gives
\begin{equation}
  1 + 2\,e^{-\pi/r^2} + 2\,e^{-4\pi/r^2} + \cdots \;\approx\; 1.842.
\end{equation}
Keeping only the leading exponential term:
$2\,e^{-\pi/r^2} \approx 0.842$, so $e^{-\pi/r^2} \approx 0.421$, giving
$\pi/r^2 \approx 0.866$ and therefore
\begin{equation}
  r \;=\; R_t / R_\psi \;\approx\; \sqrt{\pi / 0.866} \;\approx\; 1.90.
  \label{eq:r_required}
\end{equation}

A radius ratio $r \approx 1.9$ means that $R_t$ and $R_\psi$ are comparable
in magnitude.  This would imply that real time $t$ is compact at the scale
\begin{equation}
  2\pi R_t \;\sim\; 2\pi \times 1.9\,R_\psi.
\end{equation}
For $R_\psi$ of order the Planck length, $R_t$ would also be Planck-scale ---
rendering real time compact at the Planck scale.

%──────────────────────────────────────────────────────────────────────────────
\section{Lorentz Invariance Obstruction}
\label{sec:lorentz}
%──────────────────────────────────────────────────────────────────────────────

A compact real-time direction at radius $R_t$ would produce a spectrum of
KK states with masses $M_m = m/R_t$ for integers $m$.  For $m = 1$:
\begin{equation}
  M_1 \;=\; \frac{1}{R_t} \;\sim\; \frac{1}{1.9\,R_\psi}
         \;\sim\; \frac{M_\mathrm{Pl}}{1.9},
\end{equation}
where $M_\mathrm{Pl} \approx 1.22 \times 10^{19}$~GeV is the Planck mass.
These KK copies of Standard Model fields would have been observed at energies
well below $M_\mathrm{Pl}$ if $R_t \sim R_\psi$.  No such states are observed.

Furthermore, compactifying real time at any finite radius introduces
closed timelike curves (CTCs) at the classical level, which conflicts with
causality in the Standard Model sector.

\begin{deadendbox}
\textbf{Dead end (two-circle regime $r \sim 1$):}
Achieving $B_{\mathrm{eff}} \approx 46.3$ via the two-circle KK sum requires
$R_t/R_\psi \approx 1.9$, implying:
\begin{enumerate}
  \item Compact real time at the Planck scale, contradicting standard cosmology.
  \item KK partners of SM fields at $M \sim M_\mathrm{Pl}/1.9$, not observed.
  \item Closed timelike curves in the classical geometry.
\end{enumerate}
This scenario is physically inadmissible.  Approach~A3 is a dead end for
standard KK summation on $T^2$.
\end{deadendbox}

%──────────────────────────────────────────────────────────────────────────────
\section{Summary and Lesson}
\label{sec:summary}
%──────────────────────────────────────────────────────────────────────────────

\begin{center}
\begin{tabular}{lll}
  \toprule
  Regime & $B_{\mathrm{eff}}$ & Assessment \\
  \midrule
  $r = R_t/R_\psi \to \infty$ & $8\pi \approx 25.1$ (proved) & Physical; gives $B_0$ only \\
  $r = \sqrt{137} \approx 11.7$ & $8\pi + O(10^{-32})$ & Exponentially close to $B_0$ \\
  $r = \sqrt{139} \approx 11.8$ & $8\pi + O(10^{-32})$ & Exponentially close to $B_0$ \\
  $r \approx 1.9$ & $\approx 46.3$ & Physically inadmissible \\
  \bottomrule
\end{tabular}
\end{center}

\begin{gapbox}
\textbf{Verdict: DEAD END.}
The two-circle Kaluza-Klein sum on $T^2(R_t, R_\psi)$ cannot bridge the gap
between $B_0 = 8\pi$ and $B_{\mathrm{phenom}} \approx 46.3$ without
compactifying real time at the Planck scale, which is inconsistent with known
physics.  Approach~A3 is closed.

\textbf{Lesson:} The factor $1.842 = B_{\mathrm{phenom}}/B_0$ cannot arise
from a standard KK winding sum.  It must arise from a discrete arithmetic
structure (modular forms, $L$-values --- Approach A1/A4) or from the geometry
of a single quotient space (Approach A2).
\end{gapbox}

\ifdefined\INCLUDEMODE
\else
  \end{document}
\fi
